% Requires: \usepackage{booktabs, multirow}

% ─────────────────────────────────────────────
% Table 1: SVD k-value sweep (mean ± std RMSE)
% ─────────────────────────────────────────────
\begin{table}[ht]
\centering
\caption{Mean RMSE $\pm$ std.\ on test set (AIG, AmericanAirlines, APA) for each SVD
method across the three candidate values of $k$. SVD-K decomposes only the key matrix;
SVD-V only the value matrix; SVD-KV (sep.)\ decomposes both matrices independently;
SVD-KV (joint)\ decomposes both jointly.}
\label{tab:svd_k_sweep}
\begin{tabular}{l c c}
\toprule
Method & $k$ & Mean RMSE $\pm$ Std \\
\midrule
\multirow{3}{*}{SVD-K}
  & 33 & $0.2794 \pm 0.0338$ \\
  & 45 & $0.2389 \pm 0.0304$ \\
  & 62 & $\mathbf{0.1819 \pm 0.0282}$ \\
\midrule
\multirow{3}{*}{SVD-V}
  & 21 & $0.2983 \pm 0.0252$ \\
  & 45 & $0.1809 \pm 0.0173$ \\
  & 74 & $\mathbf{0.1010 \pm 0.0135}$ \\
\midrule
\multirow{3}{*}{SVD-KV (sep.)}
  & 21 & $0.4016 \pm 0.0361$ \\
  & 45 & $0.2812 \pm 0.0293$ \\
  & 68 & $\mathbf{0.1925 \pm 0.0256}$ \\
\midrule
\multirow{3}{*}{SVD-KV (joint)}
  & 37 & $0.3301 \pm 0.0339$ \\
  & 45 & $0.3076 \pm 0.0325$ \\
  & 68 & $\mathbf{0.2475 \pm 0.0300}$ \\
\bottomrule
\end{tabular}
\end{table}

\noindent Increasing $k$ monotonically reduces RMSE for every SVD variant, as expected since
more singular components retain more of the original key/value structure. SVD-V (value-only
decomposition) benefits the most from additional rank, reaching the lowest RMSE of any SVD
method ($0.1010$ at $k=74$), while SVD-KV (joint) is the weakest variant at every $k$ tested ---
jointly decomposing the concatenated $[K\,|\,V]$ matrix evidently wastes rank capacity compared
to decomposing $V$ alone. This already hints that \emph{which} matrix is decomposed matters more
than \emph{how} it is decomposed.


% ─────────────────────────────────────────────
% Table 2: Hokus Pokus epoch tuning
% ─────────────────────────────────────────────
\begin{table}[ht]
\centering
\caption{Hokus Pokus ($g_\theta$, MLP) RMSE on test set for different numbers of training
epochs ($k=45$, layer~5, head~1, 200 tokens). Bold marks the best result.
$^\dagger$Epoch at which early stopping triggered (patience of 10 epochs without improvement
in training loss); --- means no early stopping occurred within the epoch budget.}
\label{tab:hp_epochs}
\begin{tabular}{c c c}
\toprule
Epochs & Mean RMSE $\pm$ Std & ES ep.$^\dagger$ \\
\midrule
 10  & $0.2145 \pm 0.0347$ & --- \\
 20  & $0.2210 \pm 0.0400$ & --- \\
 30  & $0.2776 \pm 0.0552$ & 21  \\
 40  & $0.2129 \pm 0.0466$ & 28  \\
 50  & $0.2158 \pm 0.0468$ & 23  \\
 60  & $0.2075 \pm 0.0261$ & 22  \\
 70  & $0.1646 \pm 0.0228$ & 23  \\
 80  & $0.2241 \pm 0.0540$ & 29  \\
 90  & $\mathbf{0.1439 \pm 0.0181}$ & 31  \\
100  & $0.2418 \pm 0.0442$ & 33  \\
\bottomrule
\end{tabular}
\end{table}

\noindent RMSE does \emph{not} decrease monotonically with training length: epoch 90 gives the
best result ($0.1439$), but epochs 30, 80, and 100 are all worse than even epoch 10, and the
early-stopping epoch does not track final test RMSE either --- runs that stop earliest (epochs
60--70, ES at 22--23) are not the best performers, and the best run (epoch budget 90) stops
relatively late (ES at 31). This indicates that training $g_\theta$ is noisy and somewhat
unstable across runs --- the choice of 90 epochs used elsewhere in this report is an empirical
best-of-10 rather than a point on a smooth, predictable learning curve.


% ─────────────────────────────────────────────
% Table 3: Statistical tests — Elbow-selected k (incl. SVD Nyström, k=1)
% ─────────────────────────────────────────────
\begin{table}[H]
\centering
\caption{Pairwise Wilcoxon signed-rank tests with Bonferroni correction.
90 training epochs and $k=45$ for Hokus Pokus; SVD $k$ values (elbow-selected):
Method 1 (K)$=33$, Method 2 (V)$=21$, Method 3 (KV sep.)$=21$, Method 4 (KV joint)$=37$;
SVD Nyström $k=1$ (elbow-selected). $p_\text{adj}$ is the Bonferroni-corrected
$p$-value ($m=21$ comparisons).}
\label{tab:stat_tests_k_low}
\begin{tabular}{ll cc c}
\toprule
Method A & Method B & $p_\text{raw}$ & $p_\text{adj}$ & Reject $H_0$ \\
\midrule
Method 1 (K)       & Method 2 (V)         & $2.69 \times 10^{-8}$  & $5.65 \times 10^{-7}$  & Yes \\
Method 1 (K)       & Method 3 (KV sep.)   & $5.28 \times 10^{-14}$ & $1.11 \times 10^{-12}$ & Yes \\
Method 1 (K)       & Method 4 (KV joint)  & $5.28 \times 10^{-14}$ & $1.11 \times 10^{-12}$ & Yes \\
Method 1 (K)       & Hokus Pokus          & $5.28 \times 10^{-14}$ & $1.11 \times 10^{-12}$ & Yes \\
Method 1 (K)       & K-means              & $5.28 \times 10^{-14}$ & $1.11 \times 10^{-12}$ & Yes \\
Method 1 (K)       & SVD Nyström          & $5.28 \times 10^{-14}$ & $1.11 \times 10^{-12}$ & Yes \\
\addlinespace
Method 2 (V)       & Method 3 (KV sep.)   & $5.28 \times 10^{-14}$ & $1.11 \times 10^{-12}$ & Yes \\
Method 2 (V)       & Method 4 (KV joint)  & $4.68 \times 10^{-13}$ & $9.82 \times 10^{-12}$ & Yes \\
Method 2 (V)       & Hokus Pokus          & $5.28 \times 10^{-14}$ & $1.11 \times 10^{-12}$ & Yes \\
Method 2 (V)       & K-means              & $5.28 \times 10^{-14}$ & $1.11 \times 10^{-12}$ & Yes \\
Method 2 (V)       & SVD Nyström          & $5.28 \times 10^{-14}$ & $1.11 \times 10^{-12}$ & Yes \\
\addlinespace
Method 3 (KV sep.) & Method 4 (KV joint)  & $5.28 \times 10^{-14}$ & $1.11 \times 10^{-12}$ & Yes \\
Method 3 (KV sep.) & Hokus Pokus          & $5.28 \times 10^{-14}$ & $1.11 \times 10^{-12}$ & Yes \\
Method 3 (KV sep.) & K-means              & $5.28 \times 10^{-14}$ & $1.11 \times 10^{-12}$ & Yes \\
Method 3 (KV sep.) & SVD Nyström          & $5.28 \times 10^{-14}$ & $1.11 \times 10^{-12}$ & Yes \\
\addlinespace
Method 4 (KV joint) & Hokus Pokus         & $5.28 \times 10^{-14}$ & $1.11 \times 10^{-12}$ & Yes \\
Method 4 (KV joint) & K-means             & $5.28 \times 10^{-14}$ & $1.11 \times 10^{-12}$ & Yes \\
Method 4 (KV joint) & SVD Nyström         & $5.28 \times 10^{-14}$ & $1.11 \times 10^{-12}$ & Yes \\
\addlinespace
Hokus Pokus         & K-means             & $5.28 \times 10^{-14}$ & $1.11 \times 10^{-12}$ & Yes \\
Hokus Pokus         & SVD Nyström         & $5.28 \times 10^{-14}$ & $1.11 \times 10^{-12}$ & Yes \\
\addlinespace
K-means             & SVD Nyström         & $1.94 \times 10^{-5}$  & $4.07 \times 10^{-4}$  & Yes \\
\bottomrule
\end{tabular}
\end{table}

\noindent All 21 pairwise comparisons reject $H_0$ at this $k$ setting, so every method is
statistically distinguishable from every other. Most pairs sit at the test's numerical floor
($p_\text{raw} = 5.28\times10^{-14}$), meaning the methods are completely separated (every paired
RMSE difference has the same sign). The two exceptions are revealing: Method 1 (K) vs.\ Method 2
(V) is the least separated SVD pair ($p_\text{adj} = 5.65\times10^{-7}$), and K-means vs.\ SVD
Nyström is the least separated overall ($p_\text{adj} = 4.07\times10^{-4}$) --- with SVD Nyström
at its elbow-selected $k=1$, it becomes much closer to K-means than to any of the SVD-based methods.


% ─────────────────────────────────────────────
% Table 4: Statistical tests — k = 45 for all SVD methods (incl. SVD Nyström)
% ─────────────────────────────────────────────
\begin{table}[H]
\centering
\caption{Pairwise Wilcoxon signed-rank tests with Bonferroni correction.
90 training epochs and $k=45$ for Hokus Pokus; SVD $k=45$ for all methods (including SVD Nyström).
$p_\text{adj}$ is the Bonferroni-corrected $p$-value ($m=21$ comparisons).}
\label{tab:stat_tests_k45}
\begin{tabular}{ll cc c}
\toprule
Method A & Method B & $p_\text{raw}$ & $p_\text{adj}$ & Reject $H_0$ \\
\midrule
Method 1 (K)      & Method 2 (V)         & $5.28 \times 10^{-14}$ & $1.11 \times 10^{-12}$ & Yes \\
Method 1 (K)      & Method 3 (KV sep.)   & $5.28 \times 10^{-14}$ & $1.11 \times 10^{-12}$ & Yes \\
Method 1 (K)      & Method 4 (KV joint)  & $5.28 \times 10^{-14}$ & $1.11 \times 10^{-12}$ & Yes \\
Method 1 (K)      & Hokus Pokus          & $5.28 \times 10^{-14}$ & $1.11 \times 10^{-12}$ & Yes \\
Method 1 (K)      & K-means              & $5.28 \times 10^{-14}$ & $1.11 \times 10^{-12}$ & Yes \\
Method 1 (K)      & SVD Nyström          & $5.28 \times 10^{-14}$ & $1.11 \times 10^{-12}$ & Yes \\
\addlinespace
Method 2 (V)      & Method 3 (KV sep.)   & $5.28 \times 10^{-14}$ & $1.11 \times 10^{-12}$ & Yes \\
Method 2 (V)      & Method 4 (KV joint)  & $5.28 \times 10^{-14}$ & $1.11 \times 10^{-12}$ & Yes \\
Method 2 (V)      & Hokus Pokus          & $5.90 \times 10^{-13}$ & $1.24 \times 10^{-11}$ & Yes \\
Method 2 (V)      & K-means              & $5.28 \times 10^{-14}$ & $1.11 \times 10^{-12}$ & Yes \\
Method 2 (V)      & SVD Nyström          & $5.28 \times 10^{-14}$ & $1.11 \times 10^{-12}$ & Yes \\
\addlinespace
Method 3 (KV sep.) & Method 4 (KV joint) & $5.28 \times 10^{-14}$ & $1.11 \times 10^{-12}$ & Yes \\
Method 3 (KV sep.) & Hokus Pokus         & $5.28 \times 10^{-14}$ & $1.11 \times 10^{-12}$ & Yes \\
Method 3 (KV sep.) & K-means             & $5.28 \times 10^{-14}$ & $1.11 \times 10^{-12}$ & Yes \\
Method 3 (KV sep.) & SVD Nyström         & $5.28 \times 10^{-14}$ & $1.11 \times 10^{-12}$ & Yes \\
\addlinespace
Method 4 (KV joint) & Hokus Pokus        & $5.28 \times 10^{-14}$ & $1.11 \times 10^{-12}$ & Yes \\
Method 4 (KV joint) & K-means            & $5.28 \times 10^{-14}$ & $1.11 \times 10^{-12}$ & Yes \\
Method 4 (KV joint) & SVD Nyström        & $5.28 \times 10^{-14}$ & $1.11 \times 10^{-12}$ & Yes \\
\addlinespace
Hokus Pokus        & K-means             & $5.28 \times 10^{-14}$ & $1.11 \times 10^{-12}$ & Yes \\
Hokus Pokus        & SVD Nyström         & $5.28 \times 10^{-14}$ & $1.11 \times 10^{-12}$ & Yes \\
\addlinespace
K-means            & SVD Nyström         & $1.34 \times 10^{-11}$ & $2.81 \times 10^{-10}$ & Yes \\
\bottomrule
\end{tabular}
\end{table}

\noindent Again every pair is significant, but the pattern of which pairs are \emph{closest}
shifts: at uniform $k=45$, Method 1 (K) vs.\ Method 2 (V) and Method 2 (V) vs.\ Method 4 (KV joint)
both collapse to the numerical floor (fully separated), while the previously floor-level Method 2
(V) vs.\ Hokus Pokus pair becomes the least separated comparison in the whole table
($p_\text{adj} = 1.24\times10^{-11}$). This is the first run where an SVD method (Method 2 (V),
RMSE $0.1809$) gets statistically close to Hokus Pokus (RMSE $0.1535$) rather than being
decisively worse, foreshadowing the rank swap in the final ranking.


% ─────────────────────────────────────────────
% Table 5: Statistical tests — Best-k (95% explained variance), incl. SVD Nyström k=1
% ─────────────────────────────────────────────
\begin{table}[H]
\centering
\caption{Pairwise Wilcoxon signed-rank tests with Bonferroni correction.
90 training epochs and $k=45$ for Hokus Pokus; SVD $k$ values selected by the
95\%-explained-variance criterion: Method 1 (K)$=62$, Method 2 (V)$=74$,
Method 3 (KV sep.)$=68$, Method 4 (KV joint)$=68$; SVD Nyström $k=1$
(elbow-selected). $p_\text{adj}$ is the Bonferroni-corrected $p$-value
($m=21$ comparisons).}
\label{tab:stat_tests_best_k}
\begin{tabular}{ll cc c}
\toprule
Method A & Method B & $p_\text{raw}$ & $p_\text{adj}$ & Reject $H_0$ \\
\midrule
Method 1 (K)       & Method 2 (V)         & $5.28 \times 10^{-14}$ & $1.11 \times 10^{-12}$ & Yes \\
Method 1 (K)       & Method 3 (KV sep.)   & $1.22 \times 10^{-9}$  & $2.56 \times 10^{-8}$  & Yes \\
Method 1 (K)       & Method 4 (KV joint)  & $5.28 \times 10^{-14}$ & $1.11 \times 10^{-12}$ & Yes \\
Method 1 (K)       & Hokus Pokus          & $3.30 \times 10^{-11}$ & $6.94 \times 10^{-10}$ & Yes \\
Method 1 (K)       & K-means              & $5.28 \times 10^{-14}$ & $1.11 \times 10^{-12}$ & Yes \\
Method 1 (K)       & SVD Nyström          & $5.28 \times 10^{-14}$ & $1.11 \times 10^{-12}$ & Yes \\
\addlinespace
Method 2 (V)       & Method 3 (KV sep.)   & $5.28 \times 10^{-14}$ & $1.11 \times 10^{-12}$ & Yes \\
Method 2 (V)       & Method 4 (KV joint)  & $5.28 \times 10^{-14}$ & $1.11 \times 10^{-12}$ & Yes \\
Method 2 (V)       & Hokus Pokus          & $5.28 \times 10^{-14}$ & $1.11 \times 10^{-12}$ & Yes \\
Method 2 (V)       & K-means              & $5.28 \times 10^{-14}$ & $1.11 \times 10^{-12}$ & Yes \\
Method 2 (V)       & SVD Nyström          & $5.28 \times 10^{-14}$ & $1.11 \times 10^{-12}$ & Yes \\
\addlinespace
Method 3 (KV sep.) & Method 4 (KV joint)  & $5.28 \times 10^{-14}$ & $1.11 \times 10^{-12}$ & Yes \\
Method 3 (KV sep.) & Hokus Pokus          & $3.42 \times 10^{-13}$ & $7.19 \times 10^{-12}$ & Yes \\
Method 3 (KV sep.) & K-means              & $5.28 \times 10^{-14}$ & $1.11 \times 10^{-12}$ & Yes \\
Method 3 (KV sep.) & SVD Nyström          & $5.28 \times 10^{-14}$ & $1.11 \times 10^{-12}$ & Yes \\
\addlinespace
Method 4 (KV joint) & Hokus Pokus         & $5.28 \times 10^{-14}$ & $1.11 \times 10^{-12}$ & Yes \\
Method 4 (KV joint) & K-means             & $5.28 \times 10^{-14}$ & $1.11 \times 10^{-12}$ & Yes \\
Method 4 (KV joint) & SVD Nyström         & $5.28 \times 10^{-14}$ & $1.11 \times 10^{-12}$ & Yes \\
\addlinespace
Hokus Pokus         & K-means             & $5.28 \times 10^{-14}$ & $1.11 \times 10^{-12}$ & Yes \\
Hokus Pokus         & SVD Nyström         & $5.28 \times 10^{-14}$ & $1.11 \times 10^{-12}$ & Yes \\
\addlinespace
K-means             & SVD Nyström         & $2.67 \times 10^{-6}$  & $5.61 \times 10^{-5}$  & Yes \\
\bottomrule
\end{tabular}
\end{table}

\noindent With $k$ chosen to retain 95\% of explained variance, Method 1 (K) becomes the SVD
method with the most ambiguous separation from its neighbours: it is statistically the closest
to both Method 3 (KV sep.) ($p_\text{adj} = 2.56\times10^{-8}$) and, notably, to Hokus Pokus itself
($p_\text{adj} = 6.94\times10^{-10}$) --- the only run in which an SVD method approaches Hokus
Pokus's significance margin rather than a different SVD method doing so. Meanwhile Method 2 (V)
remains decisively separated from everything, consistent with it now achieving the single lowest
RMSE across all three runs (see \autoref{tab:ranking}).


% ─────────────────────────────────────────────
% Table 6: Combined ranking — Elbow-selected k vs. k = 45 vs. Best-k (95% var.)
% ─────────────────────────────────────────────
\begin{table}[ht]
\centering
\small
\caption{Method ranking by mean RMSE on the test set (AIG, AmericanAirlines, APA),
compared across the three runs where pairwise differences were actually tested
(Wilcoxon signed-rank with Bonferroni correction, $p_\text{adj} < 0.05$ for all
pairs in all three runs; see \autoref{tab:stat_tests_k_low}, \autoref{tab:stat_tests_k45},
and \autoref{tab:stat_tests_best_k}). \textbf{Bold} ranks indicate a method's position
changes across the three runs --- this turns out to be \emph{every} method, including
Hokus Pokus, which drops to rank 2 once the SVD methods use their 95\%-variance $k$.
Hokus Pokus (90 epochs) and K-means (8 clusters) use the same parameter in all runs.}
\label{tab:ranking}
\begin{tabular}{l ccc ccc ccc}
\toprule
& \multicolumn{3}{c}{Elbow-selected $k$} & \multicolumn{3}{c}{$k = 45$ (uniform)} & \multicolumn{3}{c}{Best-$k$ (95\% var.)} \\
\cmidrule(lr){2-4} \cmidrule(lr){5-7} \cmidrule(lr){8-10}
Method & $k$ & Rank & RMSE $\pm$ Std & $k$ & Rank & RMSE $\pm$ Std & $k$ & Rank & RMSE $\pm$ Std \\
\midrule
Hokus Pokus         & --   & \textbf{1} & $0.1540 \pm 0.0172$ & --   & \textbf{1} & $0.1535 \pm 0.0180$ & --   & \textbf{2} & $0.1541 \pm 0.0173$ \\
Method 1 (K)        & $33$ & \textbf{2} & $0.2794 \pm 0.0338$ & $45$ & \textbf{3} & $0.2389 \pm 0.0304$ & $62$ & \textbf{3} & $0.1819 \pm 0.0282$ \\
Method 2 (V)        & $21$ & \textbf{3} & $0.2983 \pm 0.0252$ & $45$ & \textbf{2} & $0.1809 \pm 0.0173$ & $74$ & \textbf{1} & $0.1010 \pm 0.0135$ \\
Method 3 (KV sep.)  & $21$ & \textbf{5} & $0.4016 \pm 0.0361$ & $45$ & \textbf{4} & $0.2812 \pm 0.0293$ & $68$ & \textbf{4} & $0.1925 \pm 0.0256$ \\
Method 4 (KV joint) & $37$ & \textbf{4} & $0.3301 \pm 0.0339$ & $45$ & \textbf{5} & $0.3076 \pm 0.0325$ & $68$ & \textbf{5} & $0.2475 \pm 0.0300$ \\
SVD Nyström         & $1$  & \textbf{6} & $0.7589 \pm 0.0263$ & $45$ & \textbf{7} & $1.2716 \pm 1.1447$ & $1$  & \textbf{6} & $0.7582 \pm 0.0269$ \\
K-means             & --   & \textbf{7} & $0.7957 \pm 0.0771$ & --   & \textbf{6} & $0.7868 \pm 0.0713$ & --   & \textbf{7} & $0.7968 \pm 0.0705$ \\
\bottomrule
\end{tabular}
\end{table}

\noindent \textbf{Main conclusions.} No method's rank is stable across the three hyperparameter
settings tested, which is itself the central finding of this comparison: \emph{any ranking of
these methods is conditional on the choice of $k$/training budget, and a ranking produced under
one setting should not be presented as the ranking.} Three points stand out:
\begin{itemize}
  \item \textbf{Hokus Pokus is the most robust choice, not the best in every case.} It is the
    only method that is never worse than rank 2 and never the best by a wide margin either ---
    its RMSE barely moves across runs ($0.1535$--$0.1541$), since its parameters (90 epochs,
    $k=45$) are fixed across all three columns. It is the safe default when the SVD rank cannot
    be tuned per-deployment.
  \item \textbf{Method 2 (V) is the best-case method, but only with the right $k$.} At $k=21$
    it is mediocre (rank 3, RMSE $0.2983$); at $k=74$ (95\%-variance-selected) it becomes the best
    method overall (RMSE $0.1010$), beating Hokus Pokus outright. Method 2 (V) has the highest
    \emph{ceiling} of any method here, but also the widest swing in rank (3 $\to$ 2 $\to$ 1),
    so realising that ceiling depends on correctly tuning $k$ for the deployment head.
  \item \textbf{K-means and SVD Nyström are reliably the two worst methods}, regardless of which
    run is consulted, and SVD Nyström is additionally unstable at higher $k$ (std balloons to
    $1.1447$ at $k=45$; see \autoref{tab:SVD_nyström}). Neither is competitive with the
    SVD-decomposition methods or Hokus Pokus at any setting tested, so they are not recommended
    candidates for this attention-approximation task.
\end{itemize}
Overall, if a single deployment-ready recommendation is required without per-head $k$-tuning,
Hokus Pokus is the preferred choice on robustness grounds; if per-head tuning of $k$ is feasible,
Method 2 (V) with a 95\%-variance-selected $k$ gives the best observed accuracy.


% ─────────────────────────────────────────────
% Table 7: SVD Nyström k sweep
% ─────────────────────────────────────────────
\noindent Results for SVD Nyström is in \autoref{tab:SVD_nyström}.
\begin{table}[H]
\centering
\begin{tabular}{l c c}
\toprule
Method & $k$ & Mean RMSE $\pm$ Std \\
\midrule
SVD Nyström & \textbf{1} & $\mathbf{0.7589 \pm 0.0263}$ \\
SVD Nyström & 45 & $1.2716 \pm 1.1447$ \\
\bottomrule
\end{tabular}
\caption{SVD Nyström RMSE on test companies (AIG, AmericanAirlines, APA).
$k=1$ (elbow-selected) gives both a lower and far more stable RMSE than $k=45$.}
\label{tab:SVD_nyström}
\end{table}

\noindent Counter-intuitively, using \emph{more} landmark points ($k=45$) makes the Nyström
approximation both worse on average ($1.2716$ vs.\ $0.7589$) and dramatically less stable
(std of $1.1447$, comparable in magnitude to the mean itself, vs.\ $0.0263$ at $k=1$). This
suggests the Nyström approximation becomes ill-conditioned as $k$ grows for this attention head
--- likely because the landmark sub-matrix approaches singularity --- rather than improving as
more rank is retained, which is the opposite of the behaviour seen for the four direct SVD
methods in \autoref{tab:svd_k_sweep}. In practice this means SVD Nyström should be used only at
very low rank for this setting, and its $k$ cannot be tuned by the same "more variance is better"
heuristic used for the other SVD methods.


% ─────────────────────────────────────────────
% Table A: Compute (GPU only) — Self-CUDA time (PyTorch profiler, layer 5, head 1, 200 tokens)
% ─────────────────────────────────────────────
\noindent\textbf{Profiling: computational cost, FLOPs, and memory.} Each method was profiled
with \texttt{torch.profiler} over 11 repeated full prefill+20-token-decode runs (layer~5,
head~1, 200-token prompt, $k=45$ / 45 clusters where applicable), all nine methods called
back-to-back on the same cached query/key/value tensors within a single profiling script
(\texttt{src/profiling/attention\_prof.py}). \autoref{tab:profiling_compute} reports both
mean~$\pm$~std and the median for \emph{self-CUDA} device time (the actual GPU kernel time
spent inside the attention call, robust to host-side scheduling noise), across all 11 runs.
The two statistics are shown together deliberately: for most methods they agree closely, but
for Method~3 (KV sep.) a single early iteration runs noticeably faster than the other ten,
which drags the mean down and inflates the std without changing the steady-state cost --- the
median is reported elsewhere in this section as the representative number for exactly this
reason, and showing both here makes that visible rather than hiding it. Host-side self-CPU
time was also profiled but is intentionally omitted from the main tables: it is dominated by
driver/JIT overhead (cuSOLVER lazy-loading, kernel-launch queuing) rather than the method's
intrinsic cost, and notably showed a $\sim\!288$\,ms anomaly for Method~1 (K) with no
counterpart in GPU time, FLOPs, or memory --- a real, reproducible measurement, but not one
that reflects computational demand in the sense this section is concerned with.
\begin{table}[ht]
\centering
\caption{Self-CUDA device time per attention call: mean~$\pm$~std and median across all 11
profiler runs (layer~5, head~1, 200 tokens). Naive and Flash are exact, uncompressed
baselines.}
\label{tab:profiling_compute}
\begin{tabular}{l c}
\toprule
Method & Self-CUDA time (mean$\pm$std / med.) \\
\midrule
Naive Attention      & \shortstack{$31.5\pm2.6\,\mu s$ \\ (med.\ $31.8$)}   \\
Flash Attention      & \shortstack{$57.5\pm2.9\,\mu s$ \\ (med.\ $58.3$)}   \\
K-means (45 clusters) & \shortstack{$53.2\pm5.3\,\mu s$ \\ (med.\ $54.9$)}  \\
Method 4 (KV joint)  & \shortstack{$4.20\pm0.12$ ms \\ (med.\ $4.22$)}      \\
Hokus Pokus          & \shortstack{$4.19\pm0.13$ ms \\ (med.\ $4.22$)}     \\
Method 1 (K)         & \shortstack{$4.19\pm0.13$ ms \\ (med.\ $4.22$)}     \\
Method 3 (KV sep.)   & \shortstack{$7.22\pm0.88$ ms$^\dagger$ \\ (med.\ $7.50$)} \\
Method 2 (V)         & \shortstack{$8.28\pm0.23$ ms \\ (med.\ $8.32$)}     \\
SVD Nyström          & \shortstack{$8.34\pm0.23$ ms \\ (med.\ $8.42$)}     \\
\bottomrule
\end{tabular}
\\[2pt]
{\footnotesize $^\dagger$Method~3 (KV sep.) shows a single-outlier pattern in GPU time: one
early iteration runs at $4.57$\,ms while the other ten cluster at $7.4$--$7.7$\,ms, pulling its
mean down to $7.22$\,ms and its std up to $0.88$\,ms (12\% of the mean, the largest relative
std of any method's CUDA time); the median of $7.50$\,ms is used elsewhere in this section as
the representative figure.}
\end{table}

\noindent Compressing the KV cache does not make attention cheaper here --- it makes it
dramatically more expensive. Every SVD-based method (Methods 1--4, SVD Nyström) and Hokus
Pokus costs \textbf{two to three orders of magnitude more GPU time per call} than the exact
Naive or Flash baselines ($31$--$58\,\mu s$), because each call re-runs a $\texttt{torch.linalg.svd}$
(or, for Hokus Pokus, an SVD plus an MLP correction) from scratch rather than reusing a
cached factorization. K-means is the only compressed method that stays close to baseline cost
($54.9\,\mu s$), since cluster assignment needs no eigendecomposition.


% ─────────────────────────────────────────────
% Table B: FLOPs (algorithmic complexity) — prefill vs. decode step
% ─────────────────────────────────────────────
\begin{table}[ht]
\centering
\caption{Floating-point operation counts per call: the prefill pass (step 0, which builds the
compressed cache) versus the mean over the following 10 single-token decode steps. Naive and
Flash are exact baselines; all other methods attend over a compressed/approximated cache.}
\label{tab:profiling_flops}
\begin{tabular}{l c c}
\toprule
Method & Prefill FLOPs & Mean decode-step FLOPs \\
\midrule
K-means (45 clusters) & $3.64$M  & $7.17$M  \\
SVD Nyström           & $7.80$M  & $14.9$M  \\
Naive Attention       & $12.8$M  & $49.7$M  \\
Flash Attention       & $12.8$M  & $49.7$M  \\
Method 1 (K)          & $15.2$M  & $54.5$M  \\
Method 4 (KV joint)   & $15.2$M  & $54.5$M  \\
Method 3 (KV sep.)    & $17.1$M  & $58.1$M  \\
Method 2 (V)          & $17.7$M  & $59.4$M  \\
Hokus Pokus           & $34.0$M  & $83.6$M  \\
\bottomrule
\end{tabular}
\end{table}

\noindent FLOPs and wall-clock GPU time \emph{disagree} for the landmark-based methods.
SVD Nyström and K-means need $3$--$7\times$ \emph{fewer} FLOPs than exact attention (they only
ever multiply against a handful of landmarks/centroids, not the full 200-token cache), yet
SVD Nyström is $\sim\!260\times$ \emph{slower} in wall-clock GPU time than Naive Attention
(\autoref{tab:profiling_compute}). The arithmetic is cheap; the $\texttt{cusolver}$ SVD kernel
launch and small-matrix decomposition overhead is not, and dominates completely at this scale.
K-means avoids this trap (low FLOPs \emph{and} low wall-clock time) precisely because it has no
decomposition step. Conversely, Methods 1--4 and Hokus Pokus need \emph{more} FLOPs than exact
attention, not fewer --- compressing the cache adds the cost of the projection/reconstruction
itself on top of attending over it, and Hokus Pokus's extra MLP correction nearly doubles its
FLOPs again relative to the plain SVD-K method it is built on (Method~1).


% ─────────────────────────────────────────────
% Table C: Memory — peak CUDA memory footprint
% ─────────────────────────────────────────────
\begin{table}[ht]
\centering
\caption{Approximate peak CUDA memory footprint per attention call (median over all 11
profiler runs), from the PyTorch profiler's reported CUDA memory column.
Measured at a 200-token context length.}
\label{tab:profiling_mem}
\begin{tabular}{l c}
\toprule
Method & Peak CUDA memory \\
\midrule
K-means (45 clusters) & $0.66$ MB \\
SVD Nyström           & $1.01$ MB \\
Flash Attention       & $2.13$ MB \\
Naive Attention       & $2.38$ MB \\
Method 1 (K)          & $2.81$ MB \\
Method 4 (KV joint)   & $2.81$ MB \\
Hokus Pokus           & $2.88$ MB \\
Method 2 (V)          & $3.24$ MB \\
Method 3 (KV sep.)    & $3.61$ MB \\
\bottomrule
\end{tabular}
\end{table}

\noindent At this short, 200-token context, compression does \emph{not} reduce memory
footprint for the SVD-based methods --- Methods 1--4 and Hokus Pokus all use \emph{more}
peak CUDA memory than either exact baseline, because the fixed overhead of materialising
SVD factors, masks, and intermediate buffers outweighs the saving from storing a rank-45
approximation instead of the full 200-token cache. Only K-means and SVD Nyström, which store
a small number of centroids/landmarks rather than any factorized reconstruction of the full
cache, actually shrink the memory footprint relative to baseline. This means the memory
benefit of cache compression is a long-context effect: it should become visible only once the
saved $O(\text{context length})$ term outgrows the fixed per-call decomposition overhead
measured here, which 200 tokens is too short to demonstrate.


% ─────────────────────────────────────────────
% Table D (optional): tradeoff summary — GPU compute, FLOPs, and memory together
% ─────────────────────────────────────────────
\noindent \textbf{Why these metrics.} Self-CUDA time and peak CUDA memory are the two columns
that actually answer "did compute cost / memory usage change," because they are the only two
metrics here that (i) measure the physical resource the question is about and (ii) are
comparable apples-to-apples across all nine methods, since every method was profiled under the
identical setup (layer~5, head~1, 200-token context, $k=45$ / 45 clusters, same profiler). One
further column is included for completeness but is secondary, for the reason below:
\begin{itemize}
  \item \textbf{Self-CUDA time} is the GPU kernel time spent on the device for that call. It
    excludes host-side scheduling and is therefore the most direct, comparable measure of
    \emph{computational cost} --- a genuine change here means the GPU did more or less work.
  \item \textbf{Peak CUDA memory} is the allocator's reported memory footprint per call, in the
    same units for every method, making it a direct measure of \emph{memory usage} with no
    proxy or conversion needed.
  \item \textbf{Decode FLOPs} is included because it is the metric compression is \emph{meant}
    to reduce, but \autoref{tab:tradeoff} below shows it does not track either Self-CUDA time
    or memory for this family of methods (SVD Nyström: fewest FLOPs, near-highest GPU time) ---
    so it cannot be used on its own to answer the question and is reported only as context.
\end{itemize}
Self-CPU (host) time is deliberately excluded from this table, for the same reason given in
\autoref{tab:profiling_compute}: it is dominated by driver/JIT overhead rather than the
method's intrinsic cost.

\begin{table}[ht]
\centering
\small
\caption{Optional summary: GPU compute cost, algorithmic cost, and memory usage together
(reconstruction loss is covered separately in \autoref{tab:ranking} and
\autoref{tab:SVD_nyström} and is intentionally omitted here; host-side CPU time is omitted for
the reason given in \autoref{tab:profiling_compute}). Self-CUDA time and peak CUDA memory are
medians over all 11 \texttt{torch.profiler} runs of the isolated attention call (layer~5,
head~1, 200-token context, $k=45$ / 45 clusters); Decode FLOPs is the mean over the 10
post-prefill decode steps, from \autoref{tab:profiling_compute}--\autoref{tab:profiling_mem}.
No end-to-end/full-model inference profiling is included in this table.}
\label{tab:tradeoff}
\begin{tabular}{l c c c}
\toprule
Method & Self-CUDA time & Decode FLOPs & Peak CUDA mem \\
\midrule
Naive Attention      & $31.8\,\mu s$ & $49.7$M & $2.38$ MB \\
Flash Attention      & $58.3\,\mu s$ & $49.7$M & $2.13$ MB \\
K-means (45 cl.)     & $54.9\,\mu s$ & $7.17$M & $0.66$ MB \\
Method 4 (KV joint)  & $4.22$ ms     & $54.5$M & $2.81$ MB \\
Hokus Pokus          & $4.22$ ms     & $83.6$M & $2.88$ MB \\
Method 1 (K)         & $4.22$ ms     & $54.5$M & $2.81$ MB \\
Method 3 (KV sep.)   & $7.50$ ms     & $58.1$M & $3.61$ MB \\
Method 2 (V)         & $8.32$ ms     & $59.4$M & $3.24$ MB \\
SVD Nyström          & $8.42$ ms     & $14.9$M & $1.01$ MB \\
\bottomrule
\end{tabular}
\end{table}

\noindent \textbf{Computational cost and memory, independent of accuracy.} Ranked purely by
GPU time, the field splits into three tiers: the two exact baselines and K-means ($32$--$58\,\mu
s$); the SVD-K-family methods and Hokus Pokus ($4.2$ ms, a $\sim\!130\times$ jump); and
Method~3, Method~2, and SVD Nyström ($7.5$--$8.4$ ms, a further $\sim\!2\times$ on top of that).
Three points are specific to cost alone, independent of how accurate each method is:
\begin{itemize}
  \item \textbf{K-means is the only compressed method that is also cheap.} It matches the exact
    baselines on GPU time, while using the least memory of any method ($0.66$ MB)
    and the fewest FLOPs ($7.17$M) --- because it never decomposes anything, only assigns tokens
    to existing centroids.
  \item \textbf{Memory usage does not track GPU time.} SVD Nyström has the highest GPU
    time ($8.42$ ms) but the second-\emph{lowest} memory footprint ($1.01$ MB); Method~3 (KV
    sep.) has the highest memory footprint ($3.61$ MB) without having the highest GPU time. Cost
    along these two axes is set by different parts of the implementation (kernel/library
    overhead for time; buffer and factor materialisation for memory) and one cannot be inferred
    from the other.
  \item \textbf{FLOPs is a poor proxy for either axis.} SVD Nyström needs the fewest FLOPs of
    any compressed method ($14.9$M, beaten only by K-means) yet has the highest GPU time;
    Hokus Pokus needs the most FLOPs ($83.6$M, from its added MLP correction) but is no slower
    in GPU time than Method~1 or Method~4 (all three at $4.22$ ms, sharing the same
    \texttt{"svd\_k"}-derived decomposition cost). FLOPs predicts neither wall-clock time nor
    memory for this family of methods.
\end{itemize}
None of the SVD-based methods (1--4, Hokus Pokus, SVD Nyström) reduce peak memory below the
exact baselines at this 200-token context; only K-means does. Every compressed method that
uses an SVD step costs at least two orders of magnitude more GPU time than exact attention,
regardless of which matrix it decomposes or how that decomposition is later used.
